\textbf{Phần 1. Trắc nghiệm lựa chọn}
\setcounterpageref{ex}{0}

\begin{ex} %[5D2V11]
Tìm chữ số $x$ thích hợp biết: $4{,}8x5 < 4{,}821$.
\choice
{$x \in \{0, 1, 2\}$}
{\True $x \in \{0, 1\}$}
{$x = 1$}
{$x = 0$}
\loigiai{
So sánh hai số thập phân có cùng phần nguyên là $4$ và cùng chữ số hàng phần mười là $8$:\\
Ở hàng phần trăm: nếu $x < 2$ thì $4{,}8x5 < 4{,}821$ luôn đúng (vì $x \in \{0, 1\}$).\\
Nếu $x = 2$, ta có $4{,}825 > 4{,}821$ (không thỏa mãn).\\
Vậy các chữ số thích hợp là $x \in \{0, 1\}$.\\
\textbf{Đáp án của câu hỏi là B}
}
\end{ex}

\begin{ex} %[5D2V11]
Có bao nhiêu số tự nhiên $x$ thỏa mãn: $14{,}28 < x < 19{,}05$?
\choice
{$4$}
{\True $5$}
{$6$}
{$7$}
\loigiai{
Vì $x$ là số tự nhiên và $14{,}28 < x < 19{,}05$ nên $x \in \{15, 16, 17, 18, 19\}$.\\
Số các số tự nhiên thỏa mãn là:
\[
19 - 15 + 1 = 5\text{ (số)}.
\]
\textbf{Đáp án của câu hỏi là B}
}
\end{ex}

\begin{ex} %[5D2V11]
Trong bốn khối kim loại có khối lượng sau: Khối A ($2\text{ kg } 450\text{ g}$), Khối B ($2{,}05\text{ kg}$), Khối C ($2\dfrac{4}{10}\text{ kg}$), Khối D ($2{,}405\text{ kg}$). Khối kim loại có khối lượng nặng nhất là
\choice
{Khối D}
{Khối C}
{\True Khối A}
{Khối B}
\loigiai{
Đổi tất cả các khối lượng về đơn vị ki-lô-gam:
\begin{itemize}
    \item Khối A: $2\text{ kg } 450\text{ g} = 2{,}45\text{ kg} = 2{,}450\text{ kg}$.
    \item Khối B: $2{,}05\text{ kg} = 2{,}050\text{ kg}$.
    \item Khối C: $2\dfrac{4}{10}\text{ kg} = 2{,}4\text{ kg} = 2{,}400\text{ kg}$.
    \item Khối D: $2{,}405\text{ kg}$.
\end{itemize}
So sánh: $2{,}050 < 2{,}400 < 2{,}405 < 2{,}450$. Khối A nặng nhất ($2{,}450\text{ kg}$).\\
\textbf{Đáp án của câu hỏi là C}
}
\end{ex}

\begin{ex} %[5D2V11]
Tìm số tự nhiên $y$ lớn nhất thỏa mãn: $y < 202{,}45$.
\choice
{$201$}
{\True $202$}
{$203$}
{$200$}
\loigiai{
Số tự nhiên lớn nhất nhỏ hơn $202{,}45$ chính là phần nguyên của số đó, tức là $202$.\\
\textbf{Đáp án của câu hỏi là B}
}
\end{ex}

\begin{ex} %[5D2V11]
Cho bốn số thập phân: $X = 3{,}125$; $Y = 3{,}152$; $Z = 3{,}1205$; $T = 3{,}1502$. Sắp xếp các số theo thứ tự giảm dần ta được
\choice
{$Y > X > T > Z$}
{\True $Y > T > X > Z$}
{$T > Y > X > Z$}
{$Y > T > Z > X$}
\loigiai{
So sánh các số:
\begin{itemize}
    \item Hàng phần trăm có chữ số $5$: $Y = 3{,}1520$ và $T = 3{,}1502 \implies Y > T$.
    \item Hàng phần trăm có chữ số $2$: $X = 3{,}1250$ và $Z = 3{,}1205 \implies X > Z$.
\end{itemize}
Thứ tự giảm dần là: $Y > T > X > Z$.\\
\textbf{Đáp án của câu hỏi là B}
}
\end{ex}

\begin{ex} %[5D2V11]
Tìm chữ số $a$ bé nhất để: $8{,}a2 > 8{,}56$.
\choice
{$a = 5$}
{\True $a = 6$}
{$a = 7$}
{$a = 4$}
\loigiai{
Nếu $a = 5$, ta có $8{,}52 < 8{,}56$ (không thỏa mãn).\\
Nếu $a \ge 6$, ta có hàng phần mười là $a > 5$ nên $8{,}a2 > 8{,}56$ luôn đúng.\\
Do đó chữ số $a$ bé nhất thỏa mãn là $a = 6$.\\
\textbf{Đáp án của câu hỏi là B}
}
\end{ex}

\begin{ex} %[5D2V11]
Bốn bạn tham gia thi nhảy xa có kết quả: Minh ($3{,}15\text{ m}$), Nam ($3{,}2\text{ m}$), Phong ($3{,}08\text{ m}$), Quân ($3{,}18\text{ m}$). Bạn giành huy chương Vàng (nhảy xa nhất) là
\choice
{Phong}
{Minh}
{Quân}
{\True Nam}
\loigiai{
Viết các số đo về cùng hai chữ số ở phần thập phân:
Minh: $3{,}15\text{ m}$; Nam: $3{,}20\text{ m}$; Phong: $3{,}08\text{ m}$; Quân: $3{,}18\text{ m}$.\\
So sánh: $3{,}08 < 3{,}15 < 3{,}18 < 3{,}20$. Nam nhảy xa nhất đạt $3{,}2\text{ m}$.\\
\textbf{Đáp án của câu hỏi là D}
}
\end{ex}

\begin{ex} %[5D2V11]
Có bao nhiêu số thập phân có một chữ số ở phần thập phân lớn hơn $7{,}5$ và nhỏ hơn $8{,}5$?
\choice
{$8$}
{\True $9$}
{$10$}
{$11$}
\loigiai{
Các số thập phân thỏa mãn là: $7{,}6; 7{,}7; 7{,}8; 7{,}9; 8{,}0; 8{,}1; 8{,}2; 8{,}3; 8{,}4$.\\
Số lượng các số là: $(84 - 76) : 1 + 1 = 9$ số.\\
\textbf{Đáp án của câu hỏi là B}
}
\end{ex}

\begin{ex} %[5D2V11]
Tìm số tự nhiên $n$ biết: $n < 5{,}87 < n + 1$.
\choice
{$n = 4$}
{\True $n = 5$}
{$n = 6$}
{$n = 7$}
\loigiai{
Vì $5 < 5{,}87 < 6$ nên $n = 5$.\\
\textbf{Đáp án của câu hỏi là B}
}
\end{ex}

\begin{ex} %[5D2V11]
Từ các thẻ số $1, 3, 5$, ghép được bao nhiêu số thập phân lớn hơn $3$ mà phần nguyên gồm một chữ số và phần thập phân gồm hai chữ số khác nhau?
\choice
{$2$}
{\True $4$}
{$6$}
{$3$}
\loigiai{
Số thập phân lớn hơn $3$ nên chữ số phần nguyên phải là $3$ hoặc $5$:
\begin{itemize}
    \item Nếu phần nguyên là $3$: hai chữ số còn lại là $1, 5$ ghép được $3{,}15$ và $3{,}51$ ($2$ số).
    \item Nếu phần nguyên là $5$: hai chữ số còn lại là $1, 3$ ghép được $5{,}13$ và $5{,}31$ ($2$ số).
\end{itemize}
Tổng cộng có: $2 + 2 = 4$ số.\\
\textbf{Đáp án của câu hỏi là B}
}
\end{ex}

\begin{ex} %[5D2V11]
Nhiệt độ cơ thể của bốn bệnh nhân đo được: Bệnh nhân A ($37{,}2^\circ\text{C}$), Bệnh nhân B ($38{,}15^\circ\text{C}$), Bệnh nhân C ($37{,}8^\circ\text{C}$), Bệnh nhân D ($38{,}05^\circ\text{C}$). Người có thân nhiệt cao nhất là
\choice
{Bệnh nhân A}
{Bệnh nhân C}
{Bệnh nhân D}
{\True Bệnh nhân B}
\loigiai{
So sánh các nhiệt độ: $37{,}20 < 37{,}80 < 38{,}05 < 38{,}15$.\\
Bệnh nhân B có thân nhiệt cao nhất ($38{,}15^\circ\text{C}$).\\
\textbf{Đáp án của câu hỏi là D}
}
\end{ex}

\begin{ex} %[5D2V11]
Cho hai số thập phân $a = 6{,}75$ và $b = 6{,}7500$. Khẳng định nào sau đây là đúng?
\choice
{$a < b$}
{$a > b$}
{\True $a = b$}
{Không thể so sánh}
\loigiai{
Khi viết thêm các chữ số $0$ vào bên phải phần thập phân của một số thập phân thì được một số thập phân bằng nó. Do đó $6{,}75 = 6{,}7500$.\\
\textbf{Đáp án của câu hỏi là C}
}
\end{ex}

\textbf{Phần 2. Trắc nghiệm đúng sai}
\setcounterpageref{ex}{0}

\begin{ex} %[5D2V11]
Cho các số thập phân: $A = 12{,}305$; $B = 12{,}35$; $C = 12{,}035$; $D = 12{,}5$. Xét tính đúng sai của các khẳng định sau:
\choiceTF
{\True Số $C$ là số nhỏ nhất trong bốn số.}
{\True Số $D$ là số lớn nhất trong bốn số.}
{So sánh giữa $A$ và $B$ ta có: $A > B$.}
{\True Viết các số theo thứ tự tăng dần là: $C < A < B < D$.}
\loigiai{
\begin{itemize}
    \item Viết thêm các chữ số $0$: $A = 12{,}305$; $B = 12{,}350$; $C = 12{,}035$; $D = 12{,}500$.
    \item So sánh: $12{,}035 < 12{,}305 < 12{,}350 < 12{,}500$.
    \item Số $C$ nhỏ nhất. Khẳng định a \textbf{Đúng}.
    \item Số $D$ lớn nhất. Khẳng định b \textbf{Đúng}.
    \item Ta có $A = 12{,}305 < 12{,}350 = B$ nên khẳng định $A > B$ là \textbf{Sai}.
    \item Thứ tự tăng dần đúng là $C < A < B < D$. Khẳng định d \textbf{Đúng}.
\end{itemize}
\begin{itemize}
    \item a) \textbf{Đúng}.
    \item b) \textbf{Đúng}.
    \item c) \textbf{Sai}.
    \item d) \textbf{Đúng}.
\end{itemize}
}
\end{ex}

\begin{ex} %[5D2V11]
Thời gian hoàn thành một bài kiểm tra môn Toán của bốn bạn học sinh:
- Lan: $32{,}4\text{ phút}$
- Hoa: $32\dfrac{1}{2}\text{ phút}$
- Mai: $32{,}25\text{ phút}$
- Cúc: $32\text{ phút } 15\text{ giây}$
Xét tính đúng sai của các khẳng định sau:
\choiceTF
{\True Thời gian của bạn Hoa là $32{,}5\text{ phút}$.}
{Bạn Lan nộp bài sớm nhất trong bốn bạn.}
{\True Đổi $32\text{ phút } 15\text{ giây} = 32{,}25\text{ phút}$.}
{\True Hai bạn Mai và Cúc nộp bài cùng một thời điểm.}
\loigiai{
\begin{itemize}
    \item Hoa: $32\dfrac{1}{2} = 32{,}5\text{ phút}$. Khẳng định a \textbf{Đúng}.
    \item Cúc: $15\text{ giây} = \dfrac{15}{60}\text{ phút} = 0{,}25\text{ phút} \implies 32{,}25\text{ phút}$. Khẳng định c \textbf{Đúng}.
    \item Mai cũng làm hết $32{,}25\text{ phút}$. Do đó Mai và Cúc làm hết thời gian bằng nhau và là hai bạn nộp sớm nhất. Khẳng định d \textbf{Đúng}.
    \item Lan làm hết $32{,}4\text{ phút} > 32{,}25\text{ phút}$ nên không phải người nộp sớm nhất. Khẳng định b \textbf{Sai}.
\end{itemize}
\begin{itemize}
    \item a) \textbf{Đúng}.
    \item b) \textbf{Sai}.
    \item c) \textbf{Đúng}.
    \item d) \textbf{Đúng}.
\end{itemize}
}
\end{ex}

\textbf{Phần 3. Trả lời ngắn}
\setcounterpageref{ex}{0}

\begin{ex} %[5D2V11]
Có bao nhiêu số tự nhiên $x$ thỏa mãn bất đẳng thức:
\[
25{,}6 < x < 31{,}2?
\]
\shortans{6}
\loigiai{
\textbf{Đáp số:} $6$\\
Các số tự nhiên thỏa mãn là: $26, 27, 28, 29, 30, 31$.\\
Số lượng các số là: $31 - 26 + 1 = 6$ (số).
}
\end{ex}

\begin{ex} %[5D2V11]
Tìm số tự nhiên $y$ lớn nhất thỏa mãn:
\[
y < 85{,}99.
\]
\shortans{85}
\loigiai{
\textbf{Đáp số:} $85$\\
Số tự nhiên lớn nhất bé hơn $85{,}99$ là $85$.
}
\end{ex}

\begin{ex} %[5D2V11]
Tìm chữ số $x$ lớn nhất để:
\[
5{,}3x8 < 5{,}372.
\]
\shortans{6}
\loigiai{
\textbf{Đáp số:} $6$\\
Hai số có cùng phần nguyên và hàng phần mười.
Nếu $x = 7$, ta có $5{,}378 > 5{,}372$ (không thỏa mãn).
Để $5{,}3x8 < 5{,}372$ thì chữ số hàng phần trăm $x$ phải nhỏ hơn $7$.
Chữ số $x$ lớn nhất thỏa mãn là $x = 6$.
}
\end{ex}

\begin{ex} %[5D2V11]
Tìm số tự nhiên $k$ bé nhất thỏa mãn:
\[
k > 99{,}005.
\]
\shortans{100}
\loigiai{
\textbf{Đáp số:} $100$\\
Số tự nhiên nhỏ nhất lớn hơn $99{,}005$ là $100$.
}
\end{ex}

\begin{ex} %[5D2V11]
Từ ba chữ số $0, 4, 9$, lập số thập phân có phần nguyên gồm một chữ số khác $0$ và phần thập phân gồm hai chữ số khác nhau. Số thập phân bé nhất lập được là bao nhiêu? (Nhập dưới dạng số thập phân, ví dụ $4.09$).
\shortans{4.09}
\loigiai{
\textbf{Đáp số:} $4.09$\\
Để số thập phân bé nhất:
Phần nguyên chọn chữ số bé nhất khác $0$ là $4$.
Hàng phần mười chọn $0$. Hàng phần trăm chọn $9$.
Số bé nhất lập được là $4{,}09$.
}
\end{ex}

\begin{ex} %[5D2V11]
Có bao nhiêu số thập phân có hai chữ số ở phần thập phân lớn hơn $3{,}15$ và bé hơn $3{,}25$?
\shortans{9}
\loigiai{
\textbf{Đáp số:} $9$\\
Các số thập phân thỏa mãn là: $3{,}16; 3{,}17; 3{,}18; 3{,}19; 3{,}20; 3{,}21; 3{,}22; 3{,}23; 3{,}24$.\\
Số các số là: $24 - 16 + 1 = 9$ số.
}
\end{ex}
